% section: 反復としての漸化式と力学系

\section{反復としての漸化式と力学系}

\subsection{数列を点の運動として見る}

一項前から次の項を決める漸化式
\[
  a_{n+1}=f(a_n)
\]
を考える.
初期値 $a_0$ を決めると,
\[
  a_0,\ f(a_0),\ f(f(a_0)),\ldots
\]
という列ができる.

この列を,関数 $f$ によって動かされる点の軌道と見る.
数学では,このように状態の更新規則を調べる分野を\textbf{力学系}という.
物理の運動を扱わなくても,値が次々に変化するなら力学系として考えられる.

この視点では,一般項よりも次の問いが中心になる.

\begin{itemize}
  \item 軌道はある値へ近づくか.
  \item ある範囲から飛び出すか.
  \item 周期的に繰り返すか.
  \item 初期値を少し変えると,軌道は大きく変わるか.
\end{itemize}

\subsection{不動点と周期点}

$f(L)=L$ を満たす点 $L$ を不動点という.
不動点から出発した軌道は,
\[
  L\longmapsto L\longmapsto L\longmapsto\cdots
\]
と動かない.

一方,$f^{\,p}(L)=L$ を満たし,$p$ 回より少ない反復では元に戻らない点を周期 $p$ の周期点という.
例えば,
\[
  f(x)=2-x
\]
では $1$ が不動点であり,$x\neq1$ ならば
\[
  x\longmapsto 2-x\longmapsto x
\]
となるので,周期 $2$ の軌道になる.

「不動点を探す」だけでは,数列の全体の動きは分からない.
不動点へ近づく軌道もあれば,そこを飛び越えて周期軌道になるものもある.

\subsection{一次関数による反復を完全に分類する}

最も基本的な反復として,
\[
  a_{n+1}=ra_n+b
\]
を考える.
$r\neq1$ のときの不動点は,
\[
  L=rL+b
\]
より,
\[
  L=\frac{b}{1-r}
\]
である.
そこで $u_n=a_n-L$ と置くと,
\[
  \begin{aligned}
    u_{n+1}
     & =a_{n+1}-L \\
     & =ra_n+b-L  \\
     & =r(a_n-L)  \\
     & =ru_n
  \end{aligned}
\]
となる.
つまり,不動点からのずれは等比数列になる.
\[
  a_n-L=r^n(a_0-L)
\]
である.

この一行から,動きが完全に分かる.

\begin{itemize}
  \item $|r|<1$ なら $a_n\to L$.
  \item $r>1$ なら,一般には不動点から遠ざかる.
  \item $r<-1$ なら,符号を交互に変えながら大きくなる.
  \item $r=-1$ なら,通常は周期 $2$ になる.
  \item $r=1$ なら,$a_{n+1}=a_n+b$ であり,等差数列になる.
\end{itemize}

ここでは,漸化式を解く方法と,反復の動きを分類する方法が一致した.
一般項を求めることが,動力学的な理解へ直結している例である.

\subsection{微分は不動点の安定性を測る}

一般の関数 $f$ について,不動点 $L$ の近くで $a_n=L+u_n$ と置く.
テイラー展開の一次近似から,
\[
  f(L+u_n)
  \approx f(L)+f'(L)u_n
  =L+f'(L)u_n
\]
となる.
したがって,
\[
  u_{n+1}\approx f'(L)u_n
\]
である.

不動点の近くでは,結局,一次関数の反復に近い.
そこで,
\[
  |f'(L)|<1
\]
ならばずれが縮み,$L$ は安定であると考えられる.
逆に,
\[
  |f'(L)|>1
\]
ならばずれが拡大し,$L$ は不安定になりやすい.

この判定は,一般項を求めなくても数列の極限を予測するための強力な道具になる.
ただし,これは $L$ の近くでの局所的な判定である.
初期値が遠い場合には,別の不動点へ行ったり,発散したりする可能性がある.

\subsection{ロジスティック写像：単純な規則から複雑な動きへ}

次の漸化式を考える.
\[
  a_{n+1}=r a_n(1-a_n)
\]
これは,人口の割合や資源の制限を簡単に表すモデルとして現れる.
$0<a_n<1$ なら,$a_n$ は割合として解釈できる.

右辺の $ra_n$ は,現在の人口に比例する増加を表す.
一方,$(1-a_n)$ は,人口が大きくなるほど増加の余地が小さくなることを表している.
一つの式の中に,増加と抑制が同時に入っている.

不動点を求めると,
\[
  L=rL(1-L)
\]
より,
\[
  L=0,\qquad L=1-\frac1r
\]
である.
後者における微分係数は,
\[
  f'(L)=r(1-2L)=2-r
\]
となる.
したがって,
\[
  |2-r|<1
\]
すなわち
\[
  1<r<3
\]
の範囲では,正の不動点が安定になると予想される.

ところが $r$ をさらに大きくすると,不動点へ収束せず,周期 $2$,周期 $4$,さらに複雑な軌道が現れる.
この現象を完全に説明するには,力学系の理論が必要になる.
しかし,ここで知っておきたいのは,非常に単純な一行の漸化式でも,一般項を求めるという問題を超えて,安定性や周期性という深い問題を生み出すということである.
